% HOW THE NULL CONTROL SEQUENCE IS SPELLED.
%
% `\csname\endcsname' makes a control sequence with NO NAME, and wherever
% the reference writes that one it writes the two control words that make
% it rather than a lone escape character:
%
%   \show      > \csname\endcsname=\relax.
%   \string    \csname\endcsname
%   context    <recently read> \csname\endcsname
%   a trace    \csname\endcsname ->X
%
% It takes the trailing space a control word takes, because that is what
% it is spelled as -- eTeX's \detokenize of it is `\csname\endcsname '
% with the space, where a one-character name would have none.
%
% A control sequence made this way is \relax to begin with, the way any
% name \csname invents is.
%
% HSTeX wrote the escape character and the empty name, which came out as
% `\' on its own. trip line 402 writes `\csname\endcsname' where an
% alignment wants a \cr, and the fault the reference draws for it names
% the sequence: `Use of \csname\endcsname doesn't match its definition.'
\catcode`\{=1 \catcode`\}=2
\tracingonline=1 \tracingmacros=2 \tracingcommands=2
\message{[A what csname invents it as]}
\expandafter\show\csname\endcsname
\message{[B defined, then shown]}
\expandafter\def\csname\endcsname{X}
\expandafter\show\csname\endcsname
\message{[C used, so the trace names it]}
\setbox0=\hbox{\csname\endcsname}
\message{[D let away, then invented again]}
\expandafter\let\csname\endcsname=\undefined
\expandafter\show\csname\endcsname
\message{[E written out as characters]}
\message{<\expandafter\string\csname\endcsname>}
\message{[F what it means]}
\message{<\expandafter\meaning\csname\endcsname>}
\message{[done]}
\end
